% ===== 原文 PDF pp.44--46（印刷 pp.43--45） =====
% tr-04 的 tr04b 片段。开头续 tr04a 末句（p.42 末：「…在一个 $A$-闭链消失的奇点附近，」）。
% 本文件 tag：3.36（印刷 p.43）；3.37--3.40（印刷 p.44）；3.41（印刷 p.45）。共 6 条。

% [原文 p.43]
则另一条闭链 $B$ 按照 Picard--Lefshetz 公式变换：
\begin{equation}
  B \to B - (B\cdot A)A
  \tag{3.36}
\end{equation}
其中 $A\cdot B$ 表示这两条闭链之间的代数相交数。

% [Tutor Note: (3.36) 按 $(B\cdot A)$ 书写；利用交截数的反对称性 $A\cdot B=-B\cdot A$，它等价于 $B\to B+(A\cdot B)A$。两种写法不是笔误，切勿擅自改符号——符号的基准是 $(A,B)$ 基的定向（$A\cdot B=\pm1$ 的取法）。]

要点在于：我们可以把求解库仑分支 EFT（即 $(a,a_D)$）的问题，转化为识别正确的 SW 描述——用 SW 曲线与微分 $(\Sigma_u,\lambda)$ 来表述——的问题，而数学上已知后者定义了一个特殊 Kähler 几何。这使我们能够以受控的方式把 $\cN=2$ 理论的非微扰物理几何化。

尚缺的是这个“涌现”环面及其微分的物理含义。这一问题的答案自然来自 6d $\cN=(2,0)$ 理论，它将是第 4 节的主要议题。

\SUB{3.4}{标度不变的 Seiberg--Witten 解与 Argyres--Douglas 理论}

我们已经简要介绍了 $\cN=2$ 理论的 Seiberg--Witten 几何 $(\Sigma_u,\lambda)$ 的一些要点；现在我们要说明，它如何为研究位于红外 RG 不动点处的 SCFT 提供线索。如前所述，非平凡的超共形不动点出现在 $\cM_{\mathrm{CB}}$ 中两个相互非局域的带电粒子同时变为无质量的那个点上；但这一点如何反映在 SW 几何 $(\Sigma_u,\lambda)$ 上呢？历史上，人们为了识别 4d $\cN=2$ 渐近自由理论的库仑分支上的 SCFT，做法是：先提取其库仑分支上的 SW 曲线与 SW 微分 $(\Sigma_u,\lambda)$；在 $\cM_{\mathrm{CB}}$ 的某些区域上，SW 几何 $(\Sigma_u,\lambda)$ 具有某种涌现的标度不变性，可用来识别超共形不动点。这正是 Argyres 与 Douglas 在 [22] 中所做的。我们可以把这种思路看作一种自顶向下（top-down）的方法，因为 SW 几何 $(\Sigma_u,\lambda)$ 是由 UV 拉格朗日量导出的。

近来，人们发展了一种更直接的方法来寻找 SCFT，即自底向上（bottom-up）的方法 [163,165,166]。它不再首先要求存在 UV 拉格朗日量，而是直接考察红外对象 $(\Sigma_u,\lambda)$，并要求该几何具有显式的标度不变性，以编码所关心理论的共形不变性。回忆一下：给定一个 4d $\cN=2$ SCFT，其库仑分支由定义 2.1 中所列手征主算子的 VEV 参数化，因而具有一个 $\mathbb{C}^*$ 作用，它由不动点处的 $\mathrm{U}(1)_r\times R^+$ 对称性诱导而来；其中 $R^+$ 指标度对称性，它被标量 VEV 自发破缺。

% [原文 p.44]
相应地，人们预期 SW 几何 $(\Sigma_u,\lambda)$（至少在局部）也具有这样的 $\mathbb{C}^*$ 作用。注意，在这种自底向上的方法中，$(\Sigma_u,\lambda)$ 这一对对象是 ad hoc 的，先验上未必与任何 $\cN=2$ 理论有关。尽管如此，事实证明，这种自底向上的方法为识别新 SCFT 提供了强有力的途径 [163,165,166]。

让我们以秩一情形为例来说明这一方法的主要思想。如前所述，Seiberg--Witten 曲线可以写成 Weierstrass 形式
\begin{equation}
  \Sigma_u:\ y^2 = x^3 + f(u)x + g(u)\,.
  \tag{3.37}
\end{equation}
我们在此选择的参数化使得 $u=0$ 时理论是标度不变的，从而是超共形的。策略是对曲线 $\Sigma_u$ 的 Weierstrass 描述施加 $\mathbb{C}^*$ 作用，以显式实现标度不变性。可以看到，这会对可能的 SW 几何施加约束。直接后果是：它要求 SW 曲线 $\Sigma_u$ 只有一个度规奇点，因为有两个奇点自然会引入一个破坏标度不变性的标度。

% [Tutor Note: 这里的 “metric singularity”（度规奇点）指的是库仑分支上有效 Kähler 度规退化的点——正是 p.39（印刷）脚注 28 所定义的奇点类型，原文在那里明确说明它不是 $\cM_{\mathrm{CB}}$ 复结构的奇点。两个这样的奇点意味着曲线上出现一个特征标度，与标度不变性冲突，故只允许一个。]

参数 $u$ 自然地带有 $\mathbb{C}^*$ 荷，因为 $u$ 应当是库仑分支上一个手征主算子的 VEV，而该算子在 $\mathrm{U}(1)_r$ R 对称性下带电。在 $\mathbb{C}^*$ 作用下，$u$ 的行为是
\begin{equation}
  u \to \xi^r u\,,
  \tag{3.38}
\end{equation}
其中 $\xi\in\mathbb{C}^*$，而 $r$ 是 $u$ 的 $\mathrm{U}(1)_r$ 荷，同时也是库仑分支手征主算子的共形维数 $\Delta$。通过在 $\Sigma_u$ 的层面上要求 $\mathbb{C}^*$ 对称性，我们得到 $f$ 与 $g$ 必须是 $u$ 的单项式。作为最后一个要素，我们利用中心荷与 BPS 粒子质量之间的关系 (3.21)，其中采用的归一化是质量在 R 对称性 $\mathrm{U}(1)_r$ 下带荷 1。同时，如式 (3.22) 所示，中心荷可以用特殊坐标表示，而特殊坐标通过式 (3.33) 与 SW 微分相联系。于是我们得出结论：SW 微分 $\lambda$ 的 $\mathrm{U}(1)_r$ 荷为
\begin{equation}
  r[\lambda] = 1\,.
  \tag{3.39}
\end{equation}
同时，由式 (3.34) 可得
\begin{equation}
  \frac{\partial \lambda}{\partial u} \sim \frac{dx}{y}\,.
  \tag{3.40}
\end{equation}
利用式 (3.39) 与 (3.40)，我们就可以得到库仑分支 $\cM_{\mathrm{CB}}$ 上手征主算子 $u$ 的标度维数：

% [Tutor Note: $f$、$g$ 必须是 $u$ 的单项式，是因为 $\mathbb{C}^*$ 作用要求曲线方程整体按一个权重缩放：若 $f$ 含 $u$ 的两个不同幂次，缩放后各项权重不一致，$u\to\xi^r u$ 就不再是曲线的一个整体缩放。这也解释了 (3.41) 的两个分母结构：单项式的指数必须为整数，故 $n$ 取正整数。]

% [原文 p.45]
\begin{equation}
  \Delta[u] = r[u] = \frac{6}{6-n}\ \text{或}\ \frac{4}{4-n} \geqslant 1\,,
  \tag{3.41}
\end{equation}
其中 $n\in\Z_+$，而它们不小于 1 的要求来自 4d 共形场论的幺正性约束\footnote{更具体地说，$d$ 维 CFT 的幺正性要求所有标量具有共形维数 $\Delta\geqslant\frac{d-2}{2}$ 或 $\Delta=0$。然而，在某些 SCFT 中，当库仑分支的坐标环不是自由生成的时候，这样的幺正性约束可以被违反：这是由于库仑分支算子之间存在非平凡关系，使得库仑分支中的坐标 $u$ 一般并不是 SCFT 中主算子的 VEV，从而它们的标度维数 $\Delta(u)$ 可以小于 1，详见 [161] 中的讨论。} [167]。事实证明，唯一可能的非平凡秩一库仑分支就是表 3 中所列的那些，其中我们略去了平凡情形 $\Delta(u)=1$。所有这些可能性都由非平凡理论实现，特别是由 Argyres--Douglas 理论实现 [22]——其定义正是：它们具有分数标度维数 $\Delta(u)$。从 Class S 的视角看，它们由 Gaiotto 曲线上的非正则穿刺点构造，我们将在 4.8.2 节介绍。

% [Tutor Note: 对 $d=4$，$\frac{d-2}{2}=1$，所以 (3.41) 中的 $\geqslant 1$ 就是 4d 标量算子的幺正性界；脚注 31 补充了“坐标环非自由生成时 $u$ 未必是主算子 VEV、$\Delta(u)$ 可小于 1”这一例外。另可逐项验算 (3.41) 与表 3：取 $6/(6-n)$ 得 $H_0$（$n=1$,$6/5$）、$H_2$（$n=2$,$3/2$）、$D_4$（$n=3$,$2$）、$\mathrm{E}_6$（$n=4$,$3$）、$\mathrm{E}_8$（$n=5$,$6$）；取 $4/(4-n)$ 得 $H_1$（$n=1$,$4/3$）与 $\mathrm{E}_7$（$n=3$,$4$）。]

这些强耦合 AD 理论的高秩版本也激发了人们通过 AdS/CFT 对应从引力角度研究它们的兴趣。最近，所有 $(A,A)$ 型 AD 理论（我们将在 4.8.2 节与第 5 节中恰当地介绍）的大中心荷情形的全息对偶已在 [168,169] 中提出。\footnote{[170] 中提出了另一个有趣的结果。作者考虑了紧化在 K3 面上的 IIB 型弦论，该 K3 面由 $n$ 个 D7 膜包裹。4d 的低能有效作用是纯 $\SU(n)$ $\cN=2$ SYM，其库仑分支上实现的是 $(A_1,A_{n-1})$ AD 理论。他们论证，每个 D7 膜分裂成一对奇异膜（exotic brane）；当 $n$ 个同类的奇异膜在同一点相撞时，这堆奇异膜的低能世界体动力学由 $(A_1,A_{n-1})$ AD 理论给出。关于这些奇异膜性质的细节，请参阅原文。}

\begin{exercisebox}{习题 3.2}
利用式 (3.39) 与 (3.40)，推导表 3 中各秩一理论里 $u$ 的共形维数，即 $\Delta[u]>1$（提示见 [63]）。
\end{exercisebox}

\SUB{3.5}{从库仑分支 EFT 读出 Argyres--Douglas 点的 SCFT 数据}

我们已经简要说明，通过施加标度不变性，SW 几何 $(\Sigma_u,\lambda)$ 如何帮助我们识别其在超共形不动点处对应的 SCFT。然而，这一自底向上的方法自始至终都存在一些注意事项。例如，SW 几何并不包含足以完全分类 4d $\cN=2$ SCFT 的信息 [163,165]。特别地，给定一条奇异的 SW 曲线及其 SW 微分，它可能对应多个不同的 SCFT。

% [Tutor Note: 这里点出了自底向上方法的根本局限：$(\Sigma_u,\lambda)$ 只编码库仑分支上的低能数据，不编码 UV 完备化；同一 SW 几何可以来自完全不同的 4d 理论（下一段 p.46 的 $\cN=4$ SYM 与 $\cN=2$ 带四个基本超多重态就是典型反例）。因此要恢复 SCFT 数据，还必须研究 SW 几何的形变（质量形变等）。]
